%% ── Section 1: Introduction ─────────────────────────────────
\section{Introduction}

Diabetic Retinopathy (DR) is a common complication of diabetes and one of the main
causes of preventable blindness in working-age adults. It damages the small blood
vessels of the retina. In fundus photographs, this damage shows up as
microaneurysms, haemorrhages and hard exudates, and in advanced cases as the growth
of new vessels. Bangladesh alone has more than 13.1~million people with
diabetes~\cite{idf2021}, and outside the big cities there are few retinal
specialists. DR is graded on a five-level scale: Grade~0 (no DR), Grade~1 (mild
non-proliferative DR, NPDR), Grade~2 (moderate NPDR), Grade~3 (severe NPDR), and
Grade~4 (proliferative DR). Referable disease (Grade~$\geq$~2) must be found early
to prevent vision loss, so automated screening is very useful.

Deep convolutional networks grade DR well when all images can be collected in one
place. On the imbalanced DDR dataset~\cite{li2019ddr}, an EfficientNet-B3 model
trained this way reaches a Quadratic Weighted Kappa (QWK) of 0.834 on the
validation split and 0.682 on the separate test split (Sec.~\ref{sec:results}).
QWK is the usual metric for ordered grades. It penalises an error by the square of
its distance, so calling a Grade~4 image Grade~0 costs much more than calling it
Grade~3. In practice, privacy rules such as HIPAA and GDPR make it hard to pool
images from several hospitals. Federated Learning (FL)~\cite{mcmahan2017fedavg}
avoids this: the hospitals train a shared model by exchanging model weights, and
the images never leave each site.

The main problem in clinical FL is that hospitals see different patients, so their
data are not identically distributed (non-IID). Screening clinics see mostly
healthy (Grade~0) eyes, referral centres see many severe cases (Grades~3--4), and
rural clinics see something in between. Under extreme Dirichlet skew
($\alpha=0.1$), standard FedAvg collapses: averaged over three seeds, its test QWK
drops from 0.647 in its best round to 0.378 in its worst. FedAvg weights each
hospital by its number of images, so a few large clinics that each see mostly one
grade can take over the global model. FedProx~\cite{li2020fedprox} avoids the
collapse by limiting how far each local model can move from the global model, but
it changes local training.

Cross-entropy, the usual training loss, also ignores the order of the grades. For
cross-entropy, mistaking Grade~4 for Grade~0 is no worse than mistaking it for
Grade~3. Ordinal losses~\cite{toledocortes2022dqor}--[\citen{polat2024cdwce}] fix
this in centralised training, but little is known about how they behave in FL.

In this paper, \textbf{OrdinalFed}, a federated method for ordinal medical grading,
is studied. It has two parts, and a repeated evaluation is used to test what each
part adds.
\begin{enumerate}[leftmargin=*]
  \item \textbf{Distance-aware OrdinalLoss:} each hospital adds a squared
  expected-grade penalty to class-weighted cross-entropy, so predictions far from
  the true grade cost more.
  \item \textbf{QWK-weighted aggregation with an adaptive $\beta$:} the server
  scores each client model on a small balanced set $\mathcal{V}$ of 50 images.
  Over the rounds, it moves the aggregation weights from image counts towards
  these QWK scores. Nothing changes at the hospitals.
  \item \textbf{A multi-seed evaluation on held-out test data:} every
  configuration is run with three seeds and scored on the DDR test split and on
  the external APTOS 2019 dataset, with ablations and a sensitivity analysis. The results show that OrdinalFed prevents
  most of FedAvg's collapse. They also show that the gain comes from giving small
  hospitals more weight, not from the QWK scores or the ordinal loss, and that
  FedProx is still the most stable method.
\end{enumerate}
